% /root/monodromy with independent suggestions and audit by /root.
% General-X A.5 proof; independently audited by /root on 19 November 2026.
\subsection{All point-count residues for arbitrary smooth ambient varieties}
\label{subsec:general-X-residues}
\status{Proved under the stated hypotheses; independently audited.}

The full generic monodromy group is unnecessary for this
consequence.  We isolate one alternating permutation factor
whose action on every other contribution has trivial trace
variation.  The argument permits a drop of the stable rank
at the reducible boundary, and includes its missing generic
unit roots in an explicit error representation.

\subsubsection{Independent node covers}

\begin{lemma}\label{lem:all-X-independent-nodes}
Let $X\subset\mathbf P^r_k$ be smooth projective of pure
dimension $n\geq2$ over an algebraically closed field.
Its connected components are denoted $X_1,\ldots,X_c$ and
have degrees $D_1,\ldots,D_c$.  Fix $e\geq3$ and $m\geq n+1$.
Take $m-1$ ordered variable ambient degree-$e$ hypersurfaces
and one of degree $e+\rho$, where $0\leq\rho<e$.
There is a nonempty open parameter space $S$ on which their
restrictions $D_i$ form a simple normal crossing divisor,
all positive-dimensional intersections on each $X_j$ are
smooth and connected, and all zero-dimensional intersections
are reduced.  On this space the joint monodromy of the
node sets
\[
 X_j\cap\bigcap_{i\in I}D_i,\qquad
 1\leq j\leq c,\quad
 I\subset\{1,\ldots,m-1\},\quad |I|=n,
\]
is the full product
\[
 \prod_{j=1}^c\prod_{|I|=n} S_{N_j},
 \qquad N_j=D_j e^n.
 \tag{6}
\]
It remains this product after restriction to any nonempty
open of $S$.
\end{lemma}

\begin{proof}
The required simultaneous smoothness and proper intersection
conditions hold on a nonempty open by Bertini, applied to
the finitely many subsets of the ordered divisors.
Positive-dimensional ample complete intersections in each
$X_j$ are geometrically connected.  This also follows by
iterating the degree-zero stable Lefschetz assertion R1.

We check the precise conditions of Entin's
Proposition~3.1, arXiv version~3, pp.~7--8
\cite{Entin2021}, for the disjoint union of these node covers
over a larger open on which they are finite.
The incidence cover for one pair $(j,I)$ is irreducible and
smooth: over $X_j$ it is a product of projective bundles
given by the linear conditions that the $n$ specified
equations vanish at the point, together with the unrestricted
other parameters.  Its fibre square off the diagonal is
irreducible for the same reason over
$X_j\times X_j\setminus\Delta$; evaluation at two distinct
points is surjective for ambient degree-$e$ forms.
Thus condition (i) of that proposition holds, and the
single-factor point action is doubly transitive.

For condition (ii), fix $(j,I)$.  Choose $n-1$ of these
divisors generally so that their intersection on $X_j$ is
a smooth connected curve $C$.  Choose a point $P\in C$.
There is a degree-$e$ equation whose restriction to $C$
has a zero of order exactly two at $P$ and only simple
zeros elsewhere, while all its intersections on the other
components are transverse.  Here are the required
jet-separation details.  In local coordinates from the
given projective embedding, the functions $1,t,t^2$ occur
as restrictions of degree-$e$ forms after multiplying by a
nonvanishing linear form to make the degrees equal.
Thus the second-order jet at $P$ can be prescribed.
The linear subsystem vanishing twice at $P$ separates
first-order jets at every other point $Q$: multiply a
degree-$(e-2)$ form by the square of a hyperplane equation
which vanishes at $P$ but not at $Q$.  Since $e-2\geq1$,
these residual forms separate tangent directions.
Bertini for this subsystem therefore gives just the one
double zero and no other nontransverse intersection.
All equations not indexed by $I$ may now be chosen generally
and not to vanish at $P$.  Every other node set stays
transverse.  The same imposed-jet Bertini argument makes
all proper partial intersections smooth.  At $P$, choose
the differential of the last equation as a linear combination
of the first $n-1$ independent differentials with every
coefficient nonzero, and prescribe its quadratic coefficient
along the tangent direction to $C$ to be nonzero.  Ambient
forms of degree at least two realize these jets.  Every
proper subtuple is then transverse at $P$, and only the
full $n$-tuple loses transversality.

The base parameter space is smooth.  The point-incidence
source is smooth at its point above $P$, hence formally
irreducible there.  Its fibre has exactly $N_j-1$ geometric
points, with all but the one double point etale.
All other node covers are etale over this parameter.
These are exactly Entin's condition (ii), valid also in
characteristic two.  Concretely, the colliding length-two
algebra has the local quadratic deformation
$z^2+az+b$; the independent constant and linear jet
parameters give the transposition even in characteristic
two.  The proposition consequently gives the full product
(6).  Removing nonempty opens preserves connectedness of
each component of its finite Galois cover, proving the
last assertion.
\end{proof}

\subsubsection{The finite part of the Lang kernel at a boundary}

\begin{lemma}\label{lem:boundary-Lang-finite-part}
Let $B$ be a smooth integral scheme over an algebraically
closed field of characteristic $p$, let $\mathcal E$ be a
vector bundle on $B$ with a Frobenius-semilinear endomorphism,
and let $U\subset B$ be a dense open on which its stable
rank is the generic rank $r$.
Let $S$ be a smooth connected affine scheme with a morphism
$f:S\to B$, such that the pulled-back stable rank is a
constant $r_0$.  After shrinking $S$ if necessary, the
special Lang local system $W$ of rank $r_0$ and any specified
finite collection of finite local systems on $S$ occur as
representations of one group $H$ which also acts on a
generic fibre $V$ of rank $r$.  There is an $H$-equivariant
inclusion $W\hookrightarrow V$, the image on the specified
finite local systems is their full joint image on $S$,
and the image on $V$ is a subgroup of the geometric
monodromy image of $\mathcal E|_U$.
\end{lemma}

\begin{proof}
Shrink to affine opens so that $f$ takes values in a smooth
affine open $B_0\subset B$.  Use the graph of $f$ in
$S\times B_0$ and henselize this smooth affine scheme along
the graph.  Write $T=\operatorname{Spec}(A)$ for the result,
with closed subscheme $\operatorname{Spec}(A/I)=S$.
The ring is Noetherian and normal, and is connected, hence integral:
henselization of a Noetherian ring is Noetherian and is a
filtered etale limit of the smooth ambient ring.  Every
connected component meets the connected closed subscheme
because $I$ is in the Jacobson radical.  The nonzero flat
map from the ambient domain to its henselization is injective,
so $T$ dominates $S\times B_0$.  The inverse image $T_U$
of $U$ is consequently a nonempty dense open.

The kernel scheme of $F-1$ on the additive vector bundle
is affine, separated, of finite presentation, and etale:
in a local trivialization its equations are
$x-Ax^{[p]}=0$, whose Jacobian is the identity.
Its fibres are finite, so it is quasi-finite.
Its restriction to $S$ is finite etale after shrinking
$S$, because the stable rank is constant.
The finite-part theorem for henselian pairs,
Stacks Lemma~41.20.5, Tag~0BTN \cite{Stacks0BTN},
splits this kernel over $T$ as a finite etale part and an
open-and-closed part disjoint from $S$.
The finite part is a subgroup: sums and inverses of its
sections stay in it. Indeed the inverse image of the away part
under multiplication or inversion is open and closed in a finite
etale scheme over the henselian pair and has no closed-fibre points;
the Jacobson-radical property makes it empty. Its restriction to
$T_U$ is a subgroup of the generic Lang local system $V$,
and its restriction to $S$ is exactly $W$.

By Stacks Lemma~58.8.2, Tag~09ZS \cite{Stacks09ZS},
finite etale covers, and their morphisms and group laws,
extend uniquely from $S$ to $T$.
Thus all the specified finite local systems extend.
The map $\pi_1(T_U)\to\pi_1(T)\simeq\pi_1(S)$ is
surjective: a connected finite etale cover of the normal
integral $T$ remains connected after restriction to a
nonempty open.  Take the joint finite image of
$\pi_1(T_U)$ on $V$ and these extended systems to be $H$.
All schemes were taken over an algebraically closed field,
so its image on $V$ lies in the desired geometric image.
No full stable-rank local system has been extended across
the rank drop; only its finite boundary part was extended.
\end{proof}

\begin{lemma}\label{lem:alternating-factor-isolation}
Let a finite group $H$ act on $\mathbf F_p$-vector spaces
$P,V,W,R,Q_0,Q_1$, with exact sequences
\[
 0\to R\to P\to M\to0,\qquad
 0\to M\to W\to Q_0\to0,\qquad
 0\to W\to V\to Q_1\to0.
 \tag{7}
\]
Suppose the image on $P$ is the product
$A=\prod_{\alpha=1}^t A_{N_\alpha}$ in its direct sum
of natural permutation representations, with every
$N_\alpha\geq\max(5,p+2)$.
Set $c=\dim R+\dim Q_0+\dim Q_1$ and
\[
 K_p=
 \begin{cases}
 2,&p=2,\\
 v_2(p^2-1)+1,&p\text{ odd}.
 \end{cases}
\]
If $t>K_pc$, then traces of elements of the image of $H$
on $V$ take every value in $\mathbf F_p$.
\end{lemma}

\begin{proof}
Let $K$ be the kernel of the action on
$R\oplus Q_0\oplus Q_1$.  Its image in the product $A$
is a normal subgroup.  Since each alternating factor is
nonabelian simple, it is the product of a subset of the
factors.  The product of the omitted factors is a quotient
of $H/K$, a subgroup of $\mathrm{GL}_c(\mathbf F_p)$.
Choose $\ell=3$ for $p=2$ and $\ell=2$ otherwise.
Each omitted factor has order divisible by $\ell$.
The number of omitted factors is therefore at most
$v_\ell(|\mathrm{GL}_c(\mathbf F_p)|)\leq K_pc$.

For completeness, the last estimate is elementary.
For $p=2$,
\[
 v_3(|\mathrm{GL}_c(\mathbf F_2)|)
 =\lfloor c/2\rfloor+
             v_3(\lfloor c/2\rfloor!)\leq2c.
\]
For odd $p$, the lifting-the-exponent formulas give
\[
 v_2(|\mathrm{GL}_c(\mathbf F_p)|)
 =c\,v_2(p-1)+\lfloor c/2\rfloor v_2(p+1)
                 +v_2(\lfloor c/2\rfloor!)
 \leq K_pc.
\]
The estimates also hold when $c=0$.
At least one factor $A_{N_\alpha}$ thus survives in the
image of $K$.  For any element $a$ of that single factor,
choose a lift $g\in K$ acting as $a$ on its permutation
summand of $P$ and identically on every other summand.
Trace additivity in (7) gives
\[
 \operatorname{Tr}(g\mid V)
  =\dim V-N_\alpha+\#\operatorname{Fix}(a)
       \quad\text{in }\mathbf F_p.
 \tag{8}
\]
Every integer $b=0,\ldots,p-1$ is the fixed-point count of
an even permutation on $N_\alpha$ letters.  Fix $b$ letters
and derange the remaining $q=N_\alpha-b\geq3$.
For odd $q$ use one $q$-cycle, and for even $q\geq4$ use
disjoint cycles of lengths $2$ and $q-2$.
These permutations are even and have exactly the required
fixed points.  Equation (8) proves the claim.
\end{proof}

\subsubsection{Explicit bounds for the error representation}

Return to a smooth projective connected
$X\subset\mathbf P^r_k$ over a finite field of characteristic
$p$, and put $L=\mathcal O_X(1)$.
The following numerical data are all computable from this
fixed embedded variety.  Cohomology dimensions are unchanged
after passage to $\bar k$ and include all its geometric
components.
Let
\[
 D=\deg_L(X),\quad A=h^{n-1}(X,\mathcal O_X),\quad
 P_\omega(t)=\chi(X,\omega_X\otimes L^t).
\]
Choose a regularity bound $r_\omega$ for the embedded
dualizing sheaf and put $d_{\rm coh}=\max(1,r_\omega-1)$.
For every $d\geq d_{\rm coh}$ the coherent relative bundle
has rank
\begin{equation}\label{eq:all-X-relative-coherent-rank}
 b_d=P_\omega(d)-h^0(X,\omega_X).
\end{equation}
Indeed Serre duality and regularity give
$H^{n-1}(X,L^{-d})=0$, and the divisor sequence identifies
the coherent cokernel with the kernel of the surjection
$H^n(X,L^{-d})\to H^n(X,\mathcal O_X)$.
This last map is onto because $H^n(Y,\mathcal O_Y)=0$.
Serre duality and the higher-cohomology vanishing for
$\omega_X\otimes L^d$ give \eqref{eq:all-X-relative-coherent-rank}.
The same construction works in families of arbitrary
effective Cartier divisors, including the boundary unions:
the dual map is multiplication by the nonzerodivisor
defining the divisor, so the displayed kernel is a vector
bundle and carries its natural semilinear Frobenius.

Choose an integer
\[
 e\geq3,\qquad e^n\geq\max(5,p+2),
 \qquad N=De^n .
 \tag{9}
\]
This guarantees that the node count on every geometric
component is at least $\max(5,p+2)$.
For $1\leq a\leq n-1$ and $t\geq0$ define the finite
cohomological upper bounds
\begin{equation}\label{eq:all-X-B-at}
 B_{a,t}=
 \sum_{j=0}^a\binom aj
 \max_{je\leq v\leq j(2e-1)}
       h^{t+j}(X,L^{-v}),
\end{equation}
where $v$ is integral and a cohomology group in degree
greater than $n$ is zero.
The Koszul resolution shows that $B_{a,t}$ bounds
$h^t(Z,\mathcal O_Z)$ for every proper complete intersection
$Z$ of $a$ divisors whose degrees lie in $[e,2e-1]$.
Its hypercohomology spectral sequence has terms
$H^{t+j}(X,L^{-\sum d_i})$ for subsets of $j$ equations,
which gives precisely the bound
\eqref{eq:all-X-B-at}.  The dimension of
$H^t_{\mathrm{et}}(Z_{\bar k},\mathbf F_p)$ is at most this
coherent dimension, by Artin--Schreier/Lang.

Set
\[
 R_0=A+\sum_{a=1}^{n-1}B_{a,n-a-1},
 \qquad
 Q_0=\sum_{a=1}^{n-1}B_{a,n-a}
           +D e^{n-1}(2e-1).
 \tag{10}
\]
For each $\rho=0,\ldots,e-1$, form the polynomial
\[
 q_\rho(z)=P_\omega(ez+\rho)-h^0(X,\omega_X)
                           -N\binom{z-1}{n}.
 \tag{11}
\]
Its degree is at most $n-1$, since the leading terms cancel.
Let $H_0$ be the maximum, over these $e$ polynomials, of
the sum of the absolute values of their coefficients.
Put
\begin{equation}\label{eq:all-X-C-bound}
\begin{aligned}
 C&=\left\lceil2R_0+Q_0+H_0\right\rceil,
       \qquad d_{\rm res}=eM_0,\\
 M_0&=\max\left\{2(n+1),\left\lceil d_{\rm coh}/e\right\rceil,
       \left\lfloor2\cdot3^{n-1}n!K_pC\right\rfloor+1\right\}.
\end{aligned}
\end{equation}
The constants can be very large; no optimality is asserted.

\begin{lemma}\label{lem:all-X-boundary-errors}
For $d=em+\rho\geq d_{\rm coh}$, with $m\geq n+1$ and
$0\leq\rho<e$, the construction of
Lemma~\ref{lem:all-X-independent-nodes} and
Lemma~\ref{lem:boundary-Lang-finite-part} gives exact
sequences of the form (7).  Here $P$ is the direct sum of
the node permutation spaces excluding the one exceptional
degree-$(e+\rho)$ component, so
$\dim P=N\binom{m-1}{n}$; $V$ is the generic relative
Lang fibre; and
\[
 \dim R+\dim Q_0^{\rm err}+\dim Q_1^{\rm err}
       \leq C(m+1)^{n-1}.
 \tag{12}
\]
The superscript distinguishes these error spaces from the
numerical constant $Q_0$ in (10).
\end{lemma}

\begin{proof}
Work first over $\bar k$ and shrink the arrangement
parameter space $S$ so that the relative stable rank is
constant and its Lang sheaf $W$ is finite locally constant.
The normal crossing cover of $Y=\bigcup_{i=1}^mD_i$
has an exact alternating resolution of its constant
$\mathbf F_p$-sheaf by the constant sheaves of its
intersections: exactness is checked on a geometric
stalk, where it is the augmented complex of a simplex.
Its cohomology spectral sequence is
\[
 E_1^{s,t}=
 \bigoplus_{|I|=s+1}
 H^t_{\mathrm{et}}(D_I,\mathbf F_p)
 \Longrightarrow H^{s+t}_{\mathrm{et}}(Y,\mathbf F_p).
 \tag{13}
\]
The last cochain term $E_1^{n-1,0}$ maps naturally to the
abutment in degree $n-1$.  Its summands from intersections
not involving the exceptional component form $P$.
Compose this map with the quotient by the ambient
$H^{n-1}_{\mathrm{et}}(X,\mathbf F_p)$ to obtain
$P\to W$, and write its kernel as $R$ and image as $M$.
This is a map of finite local systems on $S$ after the
chosen shrinking, and therefore its kernel and cokernel
are also finite local systems.

The kernel of the full last-cochain map to the abutment
is bounded by all incoming terms in the spectral sequence:
for $r=1,\ldots,n-1$ they originate at
$(n-1-r,r-1)$, of total degree $n-2$.
Taking the ambient quotient adds at most $A$ to the
kernel dimension.  Thus
\[
 \dim R\leq A+
       \sum_{a=1}^{n-1}\binom ma B_{a,n-a-1}
       \leq R_0(m+1)^{n-1}.
 \tag{14}
\]
There is no assumption that these higher differentials
vanish.

The quotient of the abutment by the full last-cochain
image is bounded by its other degree-$(n-1)$ graded terms,
namely the terms with $t>0$ in (13).  Removing from the
last-cochain space the intersections involving the
exceptional component adds at most
$D e^{n-1}(e+\rho)\binom{m-1}{n-1}$.
Passing to the relative ambient quotient cannot increase
this cokernel.  Consequently
\[
 \dim(W/M)\leq
 \sum_{a=1}^{n-1}\binom ma B_{a,n-a}
 +D e^{n-1}(2e-1)\binom{m-1}{n-1}
 \leq Q_0(m+1)^{n-1}.
 \tag{15}
\]

The finite-boundary-part lemma supplies $W\subset V$,
while $\dim V\leq b_d$.  Since
$\dim W\geq\dim P-\dim R$, we obtain
\[
 \dim(V/W)
 \leq b_d-N\binom{m-1}{n}+\dim R
 \leq(H_0+R_0)(m+1)^{n-1}.
 \tag{16}
\]
The last inequality is the coefficient bound from (11).
Equations (14)--(16) prove (12).
The exact sequences (7) now hold with
$Q_0^{\rm err}=W/M$ and $Q_1^{\rm err}=V/W$.
All node covers and these special finite systems are
included in the joint group of the boundary-part lemma;
their full independent node image is retained.
\end{proof}

\begin{theorem}\label{thm:all-X-residue-density}
Let $k$ be a finite field of characteristic $p>0$, let
$X\subset\mathbf P^r_k$ be smooth projective connected of
dimension $n\geq2$, and define the explicit integer
$d_{\rm res}$ by \eqref{eq:all-X-C-bound}.
For every $d\geq d_{\rm res}$ and every nonempty open
$\mathcal U$ in the degree-$d$ ambient hypersurface
parameter space, each $a\in\mathbf F_p$ occurs as
\[
 \#(X\cap H)(k')\equiv a\pmod p
\]
for some finite extension $k'/k$ and some
$H\in\mathcal U(k')$ with $X\cap H$ smooth.
Moreover, traces of the geometric generic relative
Hasse--Witt monodromy group take every value in
$\mathbf F_p$ in this degree range.
Geometric connectedness of $X$ is not required.
\end{theorem}

\begin{proof}
Write $d=em+\rho$ with $0\leq\rho<e$, so $m\geq M_0$.
The number of independent alternating factors is
$c_X\binom{m-1}{n}$, where $c_X\geq1$ is the number of
geometric connected components.  In particular it is
at least $\binom{m-1}{n}$.
Since $m\geq2(n+1)$,
\[
 \binom{m-1}{n}\geq\frac{m^n}{2^n n!},
 \qquad
 (m+1)^{n-1}\leq(3m/2)^{n-1}.
 \]
The strict inequality in the definition of $M_0$ therefore
gives
\[
 \binom{m-1}{n}>K_pC(m+1)^{n-1}.
 \tag{17}
\]
Restrict the joint group of the preceding lemma to the
inverse image of the product of alternating node groups.
It still surjects to that product.  Lemma
\ref{lem:alternating-factor-isolation}, with the bound (12),
shows that its traces on $V$ take every value in
$\mathbf F_p$.  This group maps into geometric monodromy
on the full generic smooth maximal-stable-rank open, so
the latter has the same trace-surjectivity.

It remains to pass from traces to the stated rational
points, with the exact extension-of-constants direction.
Replace $k$ by a finite extension splitting the geometric
components of $X$ and killing the finite Frobenius actions
on all constant spaces
$H^i_{\mathrm{et}}(X_{\bar k},\mathbf F_p)$.
We may further kill the finite arithmetic quotient of
the generic finite monodromy image.  A rational point
over any further finite extension is also a point over a
finite extension of the original $k$, so this replacement
preserves the required existential conclusion.
Every original nonempty open remains nonempty after this
base change and meets the dense smooth maximal-rank open.

Apply the finite Chebotarev theorem used earlier
\cite[Theorem~9.7.13]{KatzSarnak1999} to the finite
representation over this open.  Each geometric conjugacy
class giving a specified trace occurs over a suitable
finite extension, inside the prescribed open.
The Artin--Schreier form of weak Lefschetz R1 identifies
the lower stable cohomology of the section with that of
$X$, and identifies the remaining top quotient with $V$.
The coherent trace congruence
\cite[Corollary~3.2]{Katz1973Congruence} consequently gives
\[
 \#Y(k')\equiv
  \sum_{i=0}^{n-1}(-1)^i
       \dim_{\mathbf F_p}H^i_{\mathrm{et}}(X_{\bar k},\mathbf F_p)
  +(-1)^{n-1}\operatorname{Tr}(F_{k'}\mid V)\pmod p.
 \tag{18}
\]
Its first term is constant by the chosen field extension.
The second takes every value, which proves the claim.
If one uses arithmetic rather than geometric Frobenius
in defining the monodromy representation, replace the
chosen group element by its inverse; all conjugacy classes
are still available.  This keeps the coherent-Frobenius
trace convention in (18) unchanged.
\end{proof}

\paragraph{Scope of the conclusion.}
This proves the residue consequence A.5 with a specified
bound for every ambient variety in the question, including
the non-geometrically-connected case.  It does not identify
the full monodromy group in A.4 or the exact generic stable
rank in A.3.  In particular it never assumes generic
relative ordinarity, which the preceding smooth
$\alpha_4$ example disproves.

